% articulo-14-figuras-tablas.tex - Gráficos y figuras, pies de foto, tablas y matemáticas.
%
% FLOTANTES (Chicago fig. 3.12)
%   Objetivo: permitir que el texto continúe inmediatamente bajo una imagen.
%   \textfloatsep e \intextsep = 1 interlineado (separación flotante/texto).
%   \@fptop = 0pt: flotantes en página propia desde el borde superior.
%
% CAPTIONS
%   Separador: punto + medio cuadratín (puntocuadra).
%   Figuras: pie DEBAJO, justificado, anchura dinámica con mínimo 0.6\linewidth.
%   Tablas: título ENCIMA, cuerpo normalsize (Chicago).
%
% NOMBRES EN ESPAÑOL
%   FIGURA y TABLA en versales (convención Chicago en español).
%   Se redefinen dentro de \captionsspanish para que babel no los sobreescriba.

% ── Flotantes ─────────────────────────────────────────────────────────────────

\makeatletter
\setlength{\@fptop}{0pt plus 1fil}
\setlength{\@fpbot}{0pt plus 1fil}
\makeatother
\renewcommand{\topfraction}{1.0}
\renewcommand{\bottomfraction}{1.0}
\renewcommand{\textfraction}{0.0}
\renewcommand{\floatpagefraction}{0.5}
\setlength{\textfloatsep}{0.5\baselineskip}
\setlength{\intextsep}{0.5\baselineskip}

\usepackage{graphicx}
\usepackage{caption}
\usepackage{threeparttable}
\usepackage{tabularx}
\usepackage{multirow}

\DeclareCaptionLabelSeparator{puntocuadra}{.\enspace}

\addto\captionsspanish{%
  \renewcommand{\figurename}{FIGURA}%
  \renewcommand{\tablename}{TABLA}%
}

\captionsetup{
  font            = small,
  labelfont       = normalfont,
  labelsep        = puntocuadra,
  singlelinecheck = false,
  justification   = raggedright
}

% Anchura dinámica del caption de figura: mínimo 0.6\linewidth.
\newlength{\figcaptionwidth}
\setlength{\figcaptionwidth}{\linewidth}
\newcommand{\setfigcaptionwidth}[1]{%
  \ifdim#1>0.60\linewidth
    \setlength{\figcaptionwidth}{#1}%
  \else
    \setlength{\figcaptionwidth}{0.60\linewidth}%
  \fi
}
\captionsetup[figure]{
  position        = below,
  justification   = justified,
  singlelinecheck = false
}

% Tablas: título encima, cuerpo en normalsize.
\captionsetup[table]{
  font            = normalsize,
  position        = above,
  skip            = 4pt
}

% ── Matemáticas ───────────────────────────────────────────────────────────────

\usepackage{amsmath, amssymb}
